\documentclass[12pt]{amsart} 
\usepackage{amsfonts}
\usepackage{amscd}
\usepackage{amsmath}
\usepackage{amssymb}
\usepackage{amsthm}
\usepackage{enumitem}
\usepackage{pgf,tikz,tikz-cd}
\usepackage{mathrsfs}
\usepackage[margin=1in]{geometry}
\usepackage{array}
\usepackage{lipsum}
\usepackage{dsfont}
\usepackage{stmaryrd}

%\graphicspath{ {./diagrams/} }

\input{./latex-preamble/cm-preamble.tex}

\newtheorem*{lemma}{Lemma}
\newcommand{\G}{\mathscr{G}}
\newcommand{\sA}{\mathscr{A}}
\newcommand{\B}{\mathscr{B}}
\DeclareMathOperator{\Sh}{Sh}
\DeclareMathOperator{\PSh}{PSh}
\newcommand{\Ab}{\fr{Ab}}

\begin{document}

{Hartshorne} \hfill {1-27-2025}

\bigskip\bigskip

\begin{center}

{\sc Chapter II-1}

\end{center}

\begin{center}

  {Noah Parker} %\footnote{Collaborated with Kalev Martinson and Frank :).}}
    
\end{center}

\bigskip

% \begin{enumerate}
% 	\item[1.] {
%     }
% \end{enumerate}
% 
% \begin{proof}
% \end{proof}

\begin{enumerate}
	\item[2.] {
      \begin{enumerate}
        \item Let $\vphi:\F\to \G$ be a morphism of sheaves on $X$, and $p\in X$ a point.
          Show that $(\ker\vphi)_p\cong \ker(\vphi_p)$ and $(\im\vphi)_p\cong \im(\vphi_p)$.
        \item Show that $\vphi$ is injective (resp. surjective) if and only if $\vphi_p$ is injective (resp. surjective) for all $p\in X$.
        \item Show that a sequence $\cdots\to \F_{i-1}\xrightarrow{\vphi_{i-1}}\F_i\xrightarrow{\vphi_i}\F_{i+1}\to \cdots$ of sheaves is exact if and only if, for all $p\in X$, the associated sequence of stalks is exact.
      \end{enumerate}
    }
\end{enumerate}

\begin{proof}
  (a) A germ $(s,U)\in (\ker\vphi)_p\subset \F_p$ iff there exists $p\in V\subset U$ with $\vphi(U)(s)|_V = 0$. That is, 
  \begin{align*}
    \vphi_p(s,U) &= (\vphi(V)(s)|_V,V)=(0,V).
  \end{align*}
  Similarly, a germ $(s,U)\in (\im\vphi)_P\subset \G_p$ iff there exists $p\in V\subset U$ with $s|_V= \vphi(V)(t)$. That is,
  \begin{align*}
    (s,U) &= (s|_V,V) = (\vphi(V)(t),V) 
    = \vphi_p(t,V).
  \end{align*}

  (b) $\vphi$ is injective iff $\ker\vphi\cong 0$. The zero map $\ker\vphi\to 0$ is an isomorphism iff each of its localizations $(\ker\vphi)_p\to 0$ are isomorphisms. Using (a), we see that $\vphi$ is injective iff $\ker(\vphi_p)\cong (\ker\vphi)_p\cong 0$, i.e. $\vphi_p$ injective, for all $p\in X$.

  $\vphi$ is surjective iff $\G = \im\vphi$. That is, the inclusion map $i:\im\vphi\to \G$ is an isomorphism.
  This occurs iff each $(\im\vphi)_p\to \G_p$ is an isomorphism, which occurs iff each $\im(\vphi_p)\xrightarrow{\sim} (\im\vphi)_p\to \G_p$ is an isomorphism. 
  This map is the inclusion $\im(\vphi_p)\to \G_p$, so the result follows.

  (c) It suffices to show the result for short exact sequences $0\to \F'\xrightarrow{\vphi}\F\xrightarrow{\psi}\F''\to 0$.
  (b) gives us $\vphi$ injective iff each $\vphi_p$ is injective, $\psi$ surjective iff each $\psi_p$ is surjective. 

  If $\im\vphi=\ker\psi$, then this induces natural isomorphisms $\im\vphi_p\cong\ker\psi_p$ for all $p\in X$.
  Inclusions of sheaves induce inclusions of stalks, so $\im\vphi_p=\ker\psi_p$.

  Suppose instead that $\im\vphi_p=\ker\psi_p$ for all $p\in X$, and fix $U\subset X$, $s\in \F'(U)$.
  For each $p\in U$, we have $\wtilde{\im}\vphi_p= \ker\psi_p$, so $\psi_p\vphi_p(s,U)=(0,U)$.
  That is, there is $p\subset V_p\subset U$ such that 
  \[
    \psi(U)\vphi(U)(s)|_{V_p} =0.
  \]
  Then the $V_p$ form a cover of $U$, so the sheaf axiom gives us $\psi(U)\vphi(U)(s)=0$, whence
  \[
    \wtilde{\im}\vphi\subset \ker\psi.
  \] 
  Then, in fact, there is an inclusion $i:\im\vphi\hookrightarrow \ker\psi$ which induces the inclusions $i_p:\im\vphi_p\hookrightarrow\ker\psi_p$ for each $p\in V$.
  The latter inclusions are isomorphisms, so it follows that the former is an isomorphism as well.
  That is,
  \[
    \im\vphi=\ker\psi.
  \] 
\end{proof}

% \begin{enumerate}
% 	\item[5.] {
%       Show that a sheaf morphism $\vphi:\F\to\G$ is an isomorphism iff it is injective and surjective.
%     }
% \end{enumerate}
% 
% \begin{proof}
%   $\vphi$ is an isomorphism iff each $\vphi_p:\F_p\to \G_p$ is an isomorphism. This occurs iff each $\vphi_p$ is injective and surjective, which occurs iff $\vphi$ is injective and surjective by the 2b.
% \end{proof}

\begin{enumerate}
	\item[8.] {
      Show that, for any $U\subset X$ open, the functor $\Gamma(U,-):\Sh(X)\to \Ab$ is left-exact.
    }
\end{enumerate}

\begin{proof}
  Suppose $0\to \F'\xrightarrow{\vphi}\F\xrightarrow{\psi}\F''\to 0$ is exact.
  Then $\ker(\vphi(U)) = (\ker\vphi)(U)=0$ gives us that $\vphi(U)$ is injective. 
  It suffices to show that $\im(\vphi(U))=\ker(\psi(U))$.
  Since $\ker\psi =\im\vphi = (\wtilde{\im}\vphi)^+$, it suffices to show that $\wtilde{\im}\vphi$ is a sheaf.

  Let $\{V_i\}$ cover $V\subset X$ open, and $s_i\in \im(\vphi(V_i))$ be sections with $s_i|_{V_{ij}}=s_j|_{V_{ij}}$ for all $i,j$.
  $\vphi(V_i)$ injective gives unique $t_i\in \F'(V_i)$ such that $s_i=\vphi(V_i)(t_i)$, and
  \[
    \vphi(V_{ij})(t_i|_{V_{ij}}) = s_i|_{V_{ij}}=s_j|_{V_{ij}} = \vphi(V_{ij})(t_j|_{V_{ij}})
  \] 
  implies that the restrictions of our $t_i,t_j$ agree as well.
  Then $\F'$ a sheaf gives us a unique $t\in \F'(V)$ with $t|_{V_i}=t_i$.
  Let $s=\vphi(V)(t)\in(\wtilde{\im}\vphi)(V)$ be its image, so that $s|_{V_i} = \vphi(V_i)(t|_{V_i}) = s_i$ for each $i$.
  This shows existence.

  Suppose we have another $s'=\vphi(V)(t')$ with $s'|_{V_i}=s_i$ for each $i$. 
  Then injectivity and $\vphi(V_i)(t'|_{V_i})= \vphi(V_i)(t|_{V_i})$ implies $t'|_{V_i}=t|_{V_i}$ for all $i$, whence $t'=t$, and $s'=s$.
  We conclude that $s$ is unique, so that $\wtilde{\im}\vphi$ is a sheaf.
\end{proof}

% \begin{enumerate}
% 	\item[5.] {
%     }
% \end{enumerate}
% 
% \begin{proof}
% \end{proof}
% 
% \begin{enumerate}
%   \item[15.] {
%     }
% \end{enumerate}
% 
% \begin{proof}
% \end{proof}
% 
% \begin{enumerate}
%   \item[16.] {
%     }
% \end{enumerate}
% 
% \begin{proof}
% \end{proof}
% 
% \begin{enumerate}
%   \item[17.] {
%     }
% \end{enumerate}
% 
% \begin{proof}
% \end{proof}
% 
% \begin{enumerate}
%   \item[18.] {
%     }
% \end{enumerate}
% 
% \begin{proof}
% \end{proof}
% 
% \begin{enumerate}
%   \item[19.] {
%     }
% \end{enumerate}
% 
% \begin{proof}
% \end{proof}
% 
% \begin{enumerate}
%   \item[20.] {
%     }
% \end{enumerate}
% 
% \begin{proof}
% \end{proof}
% 
% \begin{enumerate}
%   \item[21.] {
%     }
% \end{enumerate}
% 
% \begin{proof}
% \end{proof}

\end{document}
